% Carátula del plan de tesis.
\begin{titlepage}
\centering
\vspace*{0.35cm}

{\bfseries\uppercase{Universidad Nacional de San Antonio Abad del Cusco}\par}
\vspace{0.12cm}
{\small\uppercase{Facultad de Ingeniería Eléctrica, Electrónica, Informática y Mecánica}\par}
{\small\bfseries\uppercase{Escuela Profesional de Ingeniería Informática y de Sistemas}\par}

\vspace{0.35cm}
\includegraphics[width=3.8cm]{imagenes/logo.png}\par
\vspace{0.25cm}

{\bfseries PLAN DE TESIS\par}
\vspace{0.18cm}
\rule{14.5cm}{0.8pt}\par
\vspace{0.18cm}
{\large\bfseries
\parbox{14.2cm}{\centering Diseño e implementación de una arquitectura SLM--RAG con QLoRA para la generación contextualizada en Quechua Collao}\par}
\vspace{0.18cm}
\rule{14.5cm}{0.8pt}\par

\vspace{0.28cm}
{\small Curso:}\par
{\bfseries Seminario de Investigación II\par}
\vspace{0.18cm}
{\bfseries Docente:}\par
Msc. Raúl Huillca Huallparimachi\par
\vspace{0.2cm}
{\bfseries Asesor/a de tesis:}\par
Por designar\par
\vspace{0.2cm}
{\small Para optar al Título Profesional de:}\par
{\bfseries Ingeniero Informático y de Sistemas\par}
\vspace{0.2cm}
{\bfseries Presentado por:}\par
Br. Edmil Jampier Saire Bustamante\par

\vfill
{\bfseries Cusco -- Perú\par}
{\bfseries 2026\par}
\end{titlepage}
